\documentclass[10pt,a4paper]{amsart}
\usepackage{amsmath,amssymb,amsthm,graphicx}
% For commutative diagrams we use Xy-pic
\usepackage[all]{xy}

% We use 2cell for 2-commutative diagrams.
\xyoption{2cell}
% Theorem environments.
%
\theoremstyle{plain}
\newtheorem{theorem}[subsection]{Theorem}
\newtheorem{proposition}[subsection]{Proposition}
\newtheorem{lemma}[subsection]{Lemma}

\theoremstyle{definition}
\newtheorem{definition}[subsection]{Definition}
\newtheorem{example}[subsection]{Example}
\newtheorem{exercise}[subsection]{Exercise}
\newtheorem{situation}[subsection]{Situation}

\theoremstyle{remark}
\newtheorem{remark}[subsection]{Remark}
\newtheorem{remarks}[subsection]{Remarks}

\numberwithin{equation}{subsection}

% Macros
%
\def\lim{\mathop{\mathrm{lim}}\nolimits}
\def\colim{\mathop{\mathrm{colim}}\nolimits}
\def\Spec{\mathop{\mathrm{Spec}}}
\def\Hom{\mathop{\mathrm{Hom}}\nolimits}
\def\Ext{\mathop{\mathrm{Ext}}\nolimits}
\def\SheafHom{\mathop{\mathcal{H}\!\mathit{om}}\nolimits}
\def\SheafExt{\mathop{\mathcal{E}\!\mathit{xt}}\nolimits}
\def\Sch{\mathit{Sch}}
\def\Mor{\mathop{\mathrm{Mor}}\nolimits}
\def\Ob{\mathop{\mathrm{Ob}}\nolimits}
\def\Sh{\mathop{\mathit{Sh}}\nolimits}
\def\NL{\mathop{N\!L}\nolimits}
\def\CH{\mathop{\mathrm{CH}}\nolimits}
\def\proetale{{pro\text{-}\acute{e}tale}}
\def\etale{{\acute{e}tale}}
\def\QCoh{\mathit{QCoh}}
\def\Ker{\mathop{\mathrm{Ker}}}
\def\Im{\mathop{\mathrm{Im}}}
\def\Coker{\mathop{\mathrm{Coker}}}
\def\Coim{\mathop{\mathrm{Coim}}}

% Boxtimes
%
\DeclareMathSymbol{\boxtimes}{\mathbin}{AMSa}{"02}

%
% Macros for moduli stacks/spaces
%
\def\QCohstack{\mathcal{QC}\!\mathit{oh}}
\def\Cohstack{\mathcal{C}\!\mathit{oh}}
\def\Spacesstack{\mathcal{S}\!\mathit{paces}}
\def\Quotfunctor{\mathrm{Quot}}
\def\Hilbfunctor{\mathrm{Hilb}}
\def\Curvesstack{\mathcal{C}\!\mathit{urves}}
\def\Polarizedstack{\mathcal{P}\!\mathit{olarized}}
\def\Complexesstack{\mathcal{C}\!\mathit{omplexes}}
% \Pic is the operator that assigns to X its picard group, usage \Pic(X)
% \Picardstack_{X/B} denotes the Picard stack of X over B
% \Picardfunctor_{X/B} denotes the Picard functor of X over B
\def\Pic{\mathop{\mathrm{Pic}}\nolimits}
\def\Picardstack{\mathcal{P}\!\mathit{ic}}
\def\Picardfunctor{\mathrm{Pic}}
\def\Deformationcategory{\mathcal{D}\!\mathit{ef}}
\usepackage{cleveref}
\crefname{originalreference}{original reference}{original references}
\Crefname{originalreference}{Original reference}{Original references}
\makeatletter
\newlabel{lemma-efface-cohomology-on-closed-by-finite-cover}{{09ZE}{source}{}{sourcepointer-1}{}}
\newlabel{lemma-efface-cohomology-on-closed-by-finite-cover@cref}{{[originalreference][0][]09ZE}{[1][0][]source}{}{}{}}
\makeatother
\crefname{originalreference}{original reference}{original references}
\Crefname{originalreference}{Original reference}{Original references}
\makeatletter
\newlabel{lemma-h0-henselian-pair}{{09ZH}{source}{}{sourcepointer-2}{}}
\newlabel{lemma-h0-henselian-pair@cref}{{[originalreference][0][]09ZH}{[1][0][]source}{}{}{}}
\makeatother
\title{Gabber's theorem: for a henselian pair the etale cohomology of a torsion sheaf on Spec A agrees with that on Spec A/I}
\author{The Stacks Project Authors}
\date{stacks-project:etale-cohomology.tex:local-snapshot-20261011; GNU FDL 1.2 or later}
\begin{document}
\maketitle
\noindent\textit{Editorial note.} All delivered mathematical text is the author's own, copied from the pinned Stacks Project chapter file etale-cohomology.tex; no statement or proof step was rewritten, reordered or removed.\par\medskip
\noindent\textit{Editorial note.} The counted proof is the author's complete argument: the proof is written out in full: the induction on the degree, the comparison of the two long exact sequences, and the two effacement arguments (with Z = X for injectivity and with Z = Z for surjectivity) are all present, with the degree q-1 case applied in each half.\par\medskip
\noindent\textit{Editorial note.} The label theorem-gabber names the counted statement, and the proof environment immediately following it in the source is delivered verbatim as its proof.\par\medskip
\noindent\textit{Editorial note.} The counted statement is carried by the author's own theorem environment in the chapter source.\par\medskip
\noindent\textit{Editorial note.} The statement is an author Theorem environment (named Gabber in the source), delivered verbatim with its label; the two cited inputs are kept as external pointers with their canonical Stacks tags.\par\medskip
\noindent\textit{Editorial note.} The preamble is the chapter's own native preamble excerpt already used for this source: the Xy-pic setup (native preamble lines 14-18) and the native theorem declarations with the mathematical macros (native preamble lines 154-217).\par\medskip
\noindent\textit{Allowed standard prerequisites (editorial).}\begin{itemize}
\item Henselian pair (A, I) and the closed immersion Z = Spec(A/I) into X = Spec(A). Same chapter, the notion of a henselian pair and the h\textasciicircum{}0 statement for the pair of schemes Z subset X. Supplies the geometric set-up and the vanishing lemma giving the degree zero case.
\item Torsion abelian sheaves on X\_etale and restriction along Z -> X. Same chapter, the category of torsion abelian sheaves on the etale site and the functor F |-> F|\_Z. Supplies the coefficient sheaves F, F', Q and the restricted sheaves appearing in the two long exact sequences.
\item Effacement of a cohomology class by passing to a larger torsion sheaf through a finite cover. Same chapter, the lemma effacing the cohomology of a closed subscheme by a finite cover. Supplies the choice of the injection F -> F' making an arbitrary class die in the next term, in both halves of the induction step.
\end{itemize}\medskip
\ifdefined\hypertarget\hypertarget{sourcepointer-1}{}\fi
\noindent\textit{Source reference.} Original reference 09ZE: sources/stacks-etale-cohomology/pinned-original-source.tex, source lines 14901--14901. Prerequisite pointer invoked at proof line(s) 15241; it names a standard fact of the same chapter that the counted argument uses, and every step of the counted proof built on it is retained verbatim in the delivered body.\par\medskip
\ifdefined\hypertarget\hypertarget{sourcepointer-2}{}\fi
\noindent\textit{Source reference.} Original reference 09ZH: sources/stacks-etale-cohomology/pinned-original-source.tex, source lines 15180--15180. Prerequisite pointer invoked at proof line(s) 15212; it names a standard fact of the same chapter that the counted argument uses, and every step of the counted proof built on it is retained verbatim in the delivered body.\par\medskip
\begin{theorem}[Gabber]
\label{theorem-gabber}
Let $(A, I)$ be a henselian pair. Set $X = \Spec(A)$ and
$Z = \Spec(A/I)$. For any torsion abelian sheaf $\mathcal{F}$ on $X_\etale$
we have $H^q_\etale(X, \mathcal{F}) = H^q_\etale(Z, \mathcal{F}|_Z)$.
\end{theorem}

\begin{proof}
The result holds for $q = 0$ by Lemma \ref{lemma-h0-henselian-pair}.
Let $q \geq 1$. Suppose the result has been shown in all degrees $< q$.
Let $\mathcal{F}$ be a torsion abelian sheaf. Let
$\mathcal{F} \to \mathcal{F}'$
be an injective map of torsion abelian sheaves (to be chosen later)
with cokernel $\mathcal{Q}$ so that we have the short exact sequence
$$
0 \to \mathcal{F} \to \mathcal{F}' \to \mathcal{Q} \to 0
$$
of torsion abelian sheaves on $X_\etale$. This gives a map of long exact
cohomology sequences over $X$ and $Z$ part of which looks like
$$
\xymatrix{
H^{q - 1}_\etale(X, \mathcal{F}') \ar[d] \ar[r] &
H^{q - 1}_\etale(X, \mathcal{Q}) \ar[d] \ar[r] &
H^q_\etale(X, \mathcal{F}) \ar[d] \ar[r] &
H^q_\etale(X, \mathcal{F}') \ar[d] \\
H^{q - 1}_\etale(Z, \mathcal{F}'|_Z) \ar[r] &
H^{q - 1}_\etale(Z, \mathcal{Q}|_Z) \ar[r] &
H^q_\etale(Z, \mathcal{F}|_Z) \ar[r] &
H^q_\etale(Z, \mathcal{F}'|_Z)
}
$$
Using this commutative diagram of abelian groups with exact rows
we will finish the proof.

\medskip\noindent
Injectivity for $\mathcal{F}$. Let $\xi$ be a nonzero element of
$H^q_\etale(X, \mathcal{F})$. By
Lemma \ref{lemma-efface-cohomology-on-closed-by-finite-cover} applied with
$Z = X$ (!) we can find $\mathcal{F} \subset \mathcal{F}'$ such that
$\xi$ maps to zero to the right. Then $\xi$ is the image of
an element of $H^{q - 1}_\etale(X, \mathcal{Q})$ and bijectivity
for $q - 1$ implies $\xi$ does not map to zero in
$H^q_\etale(Z, \mathcal{F}|_Z)$.

\medskip\noindent
Surjectivity for $\mathcal{F}$. Let $\xi$ be an element of
$H^q_\etale(Z, \mathcal{F}|_Z)$. By
Lemma \ref{lemma-efface-cohomology-on-closed-by-finite-cover} applied with
$Z = Z$ we can find $\mathcal{F} \subset \mathcal{F}'$ such that
$\xi$ maps to zero to the right. Then $\xi$ is the image of
an element of $H^{q - 1}_\etale(Z, \mathcal{Q}|_Z)$ and bijectivity
for $q - 1$ implies $\xi$ is in the image of the vertical map.
\end{proof}

\end{document}
